\documentclass[14pt,a4paper]{extarticle}

\usepackage[T2A]{fontenc}
\usepackage[utf8]{inputenc}
\usepackage[russian]{babel}
\usepackage{amssymb}
\usepackage{graphicx} 
\usepackage[left=3cm,right=1.5cm,top=1cm,bottom=2cm]{geometry}%размеры
\renewcommand{\baselinestretch}{1}
\renewcommand{\rmdefault}{Tempora-TLF}
\usepackage{indentfirst}
\usepackage{tikz}
\usepackage{caption}
\captionsetup{labelsep=period}
\usepackage{verbatim}
\usepackage{fancyvrb}
\usepackage{fvextra}
\usepackage{caption}
\usepackage{listings}
\usepackage{float}
\usepackage{cmap}
\usepackage[T2A]{fontenc}
\usepackage[utf8]{inputenc}
\usepackage[english,russian]{babel}
\usepackage{mathtext}
\usepackage{amsmath} 

\lstset{
    language=C++,
    basicstyle=\ttfamily\small,
    keywordstyle=\color{blue}\bfseries,
    commentstyle=\color{green},
    stringstyle=\color{red},
    numberstyle=\color{black},
    numbers=left,
    numbersep=10pt,
    columns=fullflexible,
    tabsize=4,
    showspaces=false,
    showstringspaces=true,
    breaklines=true,
    backgroundcolor=\color{white},
    frame=single,
    captionpos=b,
    escapeinside={/*@}{@*/},
    inputencoding=utf8,       
    extendedchars=\true,  
    literate={\ }{{\textcolor{white}{\textvisiblespace}}}{1},
}



\fvset{breakanywhere=true}


\begin{document}
\textcolor{blue}{\textvisiblespace}
\parindent 1cm

%титульник
{
\thispagestyle{empty}

\centering{
 Министерство образования и науки Российской Федерации\\
 \vspace{1.5mm}
 Федеральное государственное автономное образовательное учреждение высшего образования «Санкт-Петербургский политехнический университет Петра Великого»\\
 \vspace{1.5mm}
Институт компьютерных наук и кибербезопасности\\
\vspace{1.5mm}
Высшая школа технологий искусственного интелекта\\
\vspace{1.5mm}
02.03.01 Математика и компьютерные науки
 
\vspace{\fill} 
{\Large{
"Цифровая аналитика"

Отчёт о проделанной работе

\textbf{Визуальный конструктор интерфейсов.}


}}
\vspace{\fill}
    

Выполнила студентка группы 5130201/40003\hfill$\line(2,0){50} $
Рузина В.А.\\

\vspace{1.5cm}
Проверил \hfill$\line(2,0){50} $
Мулюха В.А.\\
\vspace{1cm}

\hfill<< $\line(1,0){25} $ >> $ \line(1,0){50} $ 2026 г.


\vspace{1cm}    
 \hfill г.Санкт-Петербург, 2026 г.\hfill
}}
\newpage

\section{Техническое задание}
\subsection{Цель}
{

Создание десктопного приложения для визуального проектирования веб-интерфейсов методом drag-and-drop с возможностью экспорта в чистый HTML/CSS код и настройкой всех визуальных свойств элементов.


}
{\subsection{Границы}
\subsubsection{Включено}
\begin{enumerate}
    \item Визуальный редактор с холстом для размещения элементов
    \item Палитра базовых HTML-элементов (кнопки, текстовые блоки, изображения, карточки, формы)
    \item Система drag-and-drop для добавления элементов на холст
    \item Редактор свойств для настройки внешнего вида
    \item Поддержка кастомизации CSS-свойств: цвета (текста, фона, границ), типографика (шрифт, размер, жирность, выравнивание), отступы (margin, padding), сглаживание углов (border-radius), тени (box-shadow, text-shadow), градиенты
    \item Генерация чистого HTML/CSS кода
    \item Сохранение проектов в собственном формате (.fld)
    \item Экспорт в HTML файл
    \item Масштабирование холста
    \item Выравнивание элементов (по левому краю, по центру, сетка)
\end{enumerate}

\subsubsection{Исключено}

\begin{enumerate}
    \item Редактирование HTML кода вручную
    \item Интерактивность/JavaScript на экспортируемой странице
    \item Адаптивная верстка под мобильные устройства
    \item Drag-and-drop между проектами
    \item Импорт из Figma/Photoshop
    \item Совместная работа нескольких пользователей
\end{enumerate}

\subsubsection{Ограничения на элементы}

\begin{tabular}{|p{3.5cm}|p{2.2cm}|p{3cm}|p{5.5cm}|}
\hline
\textbf{Параметр} & \textbf{Минимум} & \textbf{Максимум} & \textbf{Обоснование} \\
\hline
Количество элементов 
    & \(0\) 
    & \(10^{3}\) 
    & Нижняя граница — пустой холст; верхняя граница — производительность рендеринга \\
\hline
Размер холста (px) 
    & \(800 \times 600\) 
    & \(5000 \times 5000\) 
    & Минимум — стандартное разрешение; максимум — ограничение памяти \\
\hline
Глубина вложенности 
    & \(1\) 
    & \(10\) 
    & Избежание экспоненциального роста сложности парсинга \\
\hline
\end{tabular}

\captionof{table}{Допустимые диапазоны параметров интерфейса}

\subsection{Функционал}
{

\subsubsection {Основные компоненты интерфейса}
{
Палитра элементов

\begin{enumerate}
    \item Набор базовых блоков: кнопка, текст, изображение, карточка, контейнер, форма, поле ввода
    \item Каждый элемент имеет иконку и название
    \item Элементы перетаскиваются на холст
\end{enumerate}

Холст

\begin{enumerate}
    \item Область для размещения элементов
    \item Масштабирование
    \item Сетка для выравнивания (шаг 10px)
    \item Выделение элементов кликом
    \item Перемещение элементов мышью
    \item Изменение размера через маркеры
    \item Множественное выделение
    \item Выравнивание выделенных элементов (по левому краю, по центру, по правому краю)
\end{enumerate}

Инспектор свойств

Динамически обновляется в зависимости от выбранного элемента:

Вкладка "Стиль":

\begin{enumerate}
    \item Цвет текста
    \item Цвет фона
    \item Шрифт
    \item Размер шрифта
    \item Жирность
    \item Выравнивание текста
    \item Сглаживание углов
    \item Тень
    \item Прозрачность
\end{enumerate}

Вкладка "Геометрия":

\begin{enumerate}
    \item Позиция X, Y (в px)
    \item Ширина, Высота (в px)
    \item Отступы margin/padding (по 4 стороны)
\end{enumerate}

Вкладка "Дополнительно":

\begin{enumerate}
    \item ID элемента
    \item CSS классы
    \item Alt текст (для изображений)
    \item Ссылка (для кнопок)
\end{enumerate}
}

\subsubsection{Управление проектом}

\begin{enumerate}
    \item Создать проект - новый чистый холст
    \item Открыть проект - загрузка .fld файла
    \item Сохранить проект - сохранение в .fld
    \item Сохранить как - выбор места сохранения
    \item Экспорт в HTML - генерация HTML+CSS файла
    \item Экспорт в HTML с просмотром - экспорт и открытие в браузере
\end{enumerate}

\subsection{Технологии}

JavaScript (Node.js), Electron.js, HTML5, CSS3

\subsection{Отклонение от ТЗ}

Удалён элемент "Карточка".

Добавлен экспорт CSS-файла вместе с HTML-файлом.

Вложение элементов в контейнеры.

Управление слоями.

\newpage
\section{История промптов}

Здесь приведена хронологическая последовательность запросов в ходе разработки.

\subsection{Промпт 1 -- создание структуры проекта}

"Прочитай ТЗ из README файла и создай только структуру проекта: папки, основные файлы, README и пустые модули."

\textbf{Результат:}модель создала структуру проекта без логики. 

Было добавлено:
\begin{enumerate}
    \item package.json (line 1) с базовыми метаданными и scripts
    \item docs/architecture.md (line 1)
    \item src/main для главного процесса Electron
    \item src/preload для preload-моста
    \item src/renderer для интерфейса: компоненты палитры, холста, инспектора, toolbar
    \item src/renderer/modules с пустыми модулями под drag-and-drop, zoom, selection, export, .fld serializer
    \item src/shared для общих констант и схем
    \item tests/unit, tests/e2e
    \item .gitignore
\end{enumerate}

\subsection{Промпт 2 -- реализация логики приложения}

"Теперь реализуй интерфейс приложения согласно ТЗ. Используй HTML/CSS без сторонних UI-фреймворков."

\textbf{Результат:}модель реализовала логику и интерфейс приложения

Были внесены изменения:
\begin{enumerate}
    \item index.html (line 1): верхняя панель, палитра, холст, инспектор свойств.
    \item renderer.js (line 1): добавление элементов, drag-and-drop, выделение, множественное выделение через Shift, перемещение, resize-маркеры, zoom, сетка 10px, выравнивание, сохранение .fld, открытие .fld, экспорт HTML, предпросмотр.
    \item base.css (line 1), layout.css (line 1), canvas.css (line 1): полноценная рабочая компоновка приложения.
    \item src/main/index.js (line 1): минимальный запуск Electron-окна.
    \item src/preload/index.js (line 1): базовый preload.
\end{enumerate}

\subsection{Промпт 3 -- внесение правок}

"Нужно сделать множественное выделение мышкой

Кнопка должна быть как элемент кнопка, а не ссылка

Нужно добавить управление слоями элементами

Добавить удаление элемента с холста

В контейнере сделать возможность добавление элемента в него

В элементе текста учитывать ввод на перенос на новую строку

В элементе "Форма" нужна возможность добавления строк ввода"

\textbf{Результат:}модель исправила все недочёты, которые были описаны в промпте и изменила файлы renderer.js (line 1), index.html (line 1), layout.css (line 1) и canvas.css (line 1)

\subsection{Промпт 4 -- внесение правок для элемента "Форма" и экспорта HTML и CSS файлов}

"В форме должна быть кнопка "Отправить форму", и у этой кнопки можно менять текст, форму, шрифт, дизайн и тд
При экспорте сохраняй именно папку, где лежит html и css отдельными файлами

При предпросмотре форма меньшего размера, чем на холсте и из-за этого не видны все поля ввода, нужно исправить
Во вкладке "Дополнительно" пункт "Alt текст" должно быть только для изображения, а пункт "Ссылка" только для кнопок"

\textbf{Результат:}модель исправила все недочёты, которые были описаны в промпте. Но для элемента "Форма" были исправлены не все недочёты, появилась проблема с неправильным отображением в предпросмотре страницы, которую нужно было исправить.

\subsection{Промпт 5 -- доисправление элемента "Форма" и удаление элемента "Карточка"}

"Удали элемент "Карточка"

Элемент "Форма" всё равно на предпросмотре не соответствует размеру как на холсте

Также в элементе "Форма" нужно сделать минимальный размер элемента, чтобы нельзя было уменьшить размер так, что не влезали бы все строки на ввод и кнопку"

\textbf{Результат:}одель исправила все недочёты, которые были описаны в промпте.

\subsection{Промпт 6 -- исправление бага}

"Исправь баг:
Когда редактирую текст одного элемента и сразу переключаюсь на другой, то автоматически перекидывается изменение и на второй элемент"

\textbf{Результат:}модель исправила баг согласно промпту. Модель исправила баг в файле renderer.js

\newpage
\section{Описание использованной языковой модели}

\subsection{Семейство и версия}

В качестве инструмента разработки использовалась языковая модель Codex производства OpenAI в режиме интерактивного агента Codex в среде Visual Studio Code. Codex является специализированным агентом для работы с программным кодом: он может читать файлы проекта, изменять исходный код, запускать команды и помогать при исправлении ошибок.

\subsection{Технические характеристики}

\begin{itemize}
    \item Семейство: Codex --- специализированная линейка OpenAI для задач программирования.

    \item Режим работы: интерактивный агент в IDE, работающий непосредственно с файлами проекта.

    \item Поддерживаемые действия: чтение файлов, редактирование кода, создание новых файлов, запуск команд и анализ ошибок сборки.

    \item Среда применения: Visual Studio Code с установленным расширением Codex.

    \item Назначение: генерация исходного кода, исправление ошибок, анализ структуры проекта и помощь в реализации программного средства.

    \item Модель: в Codex может использоваться специализированная модель GPT-5.1-Codex, оптимизированная для длительных агентных задач программирования.
\end{itemize}

\subsection{Особенности применения}

В рамках данной работы модель использовалась как:

\begin{enumerate}
    \item Архитектор --- проектирование структуры приложения, разбиение проекта на модули и определение основных компонентов.

    \item Исполнитель --- генерация исходных файлов проекта, реализация интерфейса, логики drag-and-drop, инспектора свойств, сохранения проекта и экспорта в HTML/CSS.

    \item Отладчик --- корректировка исходного кода.

    \item Технический помощник --- пояснений к структуре проекта.
\end{enumerate}

\newpage
\section{Архитектура реализованного приложения}

Проект реализован как Electron-приложение с основной логикой визуального конструктора в renderer-слое. Приложение состоит из нескольких логических частей.

\begin{enumerate}
    \item \textbf{Главный процесс Electron (\texttt{src/main/})}:

    Основной файл \texttt{src/main/index.js} отвечает за desktop-оболочку приложения:
    \begin{itemize}
        \item создаёт главное окно приложения через \texttt{BrowserWindow};
        \item загружает интерфейс из файла \texttt{src/renderer/index.html};
        \item подключает preload-скрипт;
        \item обрабатывает IPC-команду \texttt{export-project-folder};
        \item открывает системный диалог выбора папки для экспорта;
        \item записывает сгенерированные файлы \texttt{index.html} и \texttt{styles.css} в выбранную папку.
    \end{itemize}

    \item \textbf{Preload-слой (\texttt{src/preload/})}:

    Файл \texttt{src/preload/index.js} используется как безопасный мост между главным процессом Electron и пользовательским интерфейсом:
    \begin{itemize}
        \item предоставляет объект \texttt{window.builderApi};
        \item передаёт наружу информацию о платформе;
        \item содержит метод \texttt{exportProject(payload)}, который вызывает IPC-команду в main-процессе.
    \end{itemize}

    Благодаря этому renderer-слой не получает прямого доступа к Node.js, что повышает безопасность Electron-приложения.

    \item \textbf{Renderer-слой (\texttt{src/renderer/})}:

    Основной файл \texttt{src/renderer/renderer.js} содержит главную логику визуального редактора. В нём реализованы:
    \begin{itemize}
        \item список элементов палитры;
        \item стандартные настройки элементов;
        \item состояние проекта;
        \item отрисовка палитры, холста и инспектора свойств;
        \item механизм drag-and-drop;
        \item выделение и множественное выделение элементов;
        \item изменение размеров элементов;
        \item вложение элементов в контейнеры;
        \item управление слоями;
        \item сохранение проекта в формате \texttt{.fld};
        \item загрузка проекта из файла \texttt{.fld};
        \item экспорт проекта в HTML/CSS;
        \item предпросмотр результата.
    \end{itemize}

    Архитектурно renderer-слой пока является монолитным модулем, однако внутри файла уже выделены отдельные зоны ответственности.

    \item \textbf{Модель состояния проекта}:

    Основное состояние приложения хранится в объекте \texttt{state}:

    \begin{verbatim}
{
  name,
  canvas,
  elements,
  selectedIds,
  zoom,
  grid,
  activeTab
}
    \end{verbatim}

    Каждый элемент в массиве \texttt{state.elements} имеет следующую структуру:

    \begin{verbatim}
{
  id,
  type,
  parentId,
  x,
  y,
  width,
  height,
  text,
  alt,
  src,
  href,
  classes,
  formFields,
  submitButton,
  style
}
    \end{verbatim}

    Вложенность элементов реализована через поле \texttt{parentId}. Если элемент находится внутри контейнера, то его \texttt{parentId} равен идентификатору этого контейнера.

    \item \textbf{Пользовательский интерфейс}:

    Разметка приложения находится в файле \texttt{src/renderer/index.html}. Интерфейс состоит из трёх основных зон:
    \begin{itemize}
        \item палитра элементов слева;
        \item рабочий холст по центру;
        \item инспектор свойств справа.
    \end{itemize}

    Стили приложения разделены на несколько файлов:
    \begin{itemize}
        \item \texttt{base.css} --- базовые переменные, сброс стандартных стилей и общие правила;
        \item \texttt{layout.css} --- панели, верхняя панель, инспектор свойств и поля ввода;
        \item \texttt{canvas.css} --- холст, элементы на холсте, маркеры изменения размера, формы и рамки выделения.
    \end{itemize}

    \item \textbf{Экспорт проекта}:

    Экспорт реализован в два этапа. На первом этапе renderer-слой генерирует:
    \begin{itemize}
        \item HTML-код через функцию \texttt{exportHtml()};
        \item CSS-код через функцию \texttt{exportStyles()}.
    \end{itemize}

    На втором этапе сгенерированные данные передаются в главный процесс через метод \texttt{window.builderApi.exportProject}. Главный процесс создаёт папку экспорта и записывает в неё файлы \texttt{index.html} и \texttt{styles.css}. Для браузерного fallback-режима предусмотрено скачивание файлов отдельно.

    \item \textbf{Сохранение и загрузка проекта}:

    Формат \texttt{.fld} представляет собой JSON-сериализацию текущего состояния проекта:

    \begin{verbatim}
{
  version: 1,
  name,
  canvas,
  elements
}
    \end{verbatim}

    При загрузке файла \texttt{.fld} элементы проходят нормализацию через функцию \texttt{normalizeElement()}. Это позволяет избежать ошибок при открытии старых сохранений после добавления новых свойств элементов.

    \item \textbf{Документация и тесты}:

    Папка \texttt{docs} содержит документацию по архитектуре проекта. Папка \texttt{tests} используется как заготовка для будущих тестов приложения.
\end{enumerate}

В текущей версии проекта часть модульных файлов уже подготовлена, но основная логика пока находится в \texttt{renderer.js}. К таким файлам относятся:

\begin{itemize}
    \item \texttt{selection.js} --- логика выделения элементов;
    \item \texttt{resize-handles.js} --- логика изменения размеров;
    \item \texttt{project-store.js} --- хранение состояния проекта и создание элементов;
    \item \texttt{html-exporter.js} --- генерация HTML/CSS;
    \item \texttt{fld-serializer.js} --- сохранение и загрузка файлов \texttt{.fld}.
\end{itemize}

\newpage
\section{Заключение и выводы}

\subsection{Итоги}

В ходе работы было разработано полнофункциональное desktop-приложение для визуального проектирования веб-интерфейсов методом drag-and-drop, удовлетворяющее требованиям технического задания. Реализован визуальный редактор с рабочим холстом, палитрой элементов, инспектором свойств и системой экспорта проектов в HTML/CSS.

Дополнительно к основным требованиям были реализованы:
\begin{itemize}
    \item вложение элементов в контейнеры;
    \item управление слоями;
    \item экспорт проекта в отдельные файлы \texttt{index.html} и \texttt{styles.css};
\end{itemize}

Архитектура приложения построена на Electron и разделена на main-process, preload-слой и renderer-слой, что обеспечивает разделение ответственности между пользовательским интерфейсом и системными функциями приложения.

\subsection{Достоинства}

\begin{itemize}
    \item Реализован полноценный визуальный редактор интерфейсов с drag-and-drop взаимодействием, позволяющий создавать пользовательские интерфейсы без ручного написания HTML/CSS-кода.

    \item Приложение использует архитектуру Electron с разделением на main-process, preload и renderer, что соответствует современному подходу к построению desktop-приложений на веб-технологиях.

    \item Реализована безопасная схема взаимодействия renderer-слоя с Node.js через preload API, что повышает безопасность приложения.

    \item Формат \texttt{.fld} основан на JSON-сериализации состояния проекта, благодаря чему структура проекта легко сохраняется и восстанавливается.

    \item Система экспорта генерирует чистый HTML/CSS, пригодный для дальнейшего использования вне приложения.

    \item Интерфейс разделён на независимые зоны: палитру элементов, холст и инспектор свойств, что упрощает дальнейшее расширение функционала.

    \item Использование Electron, HTML5, CSS3 и JavaScript обеспечивает кроссплатформенность приложения.

    \item Архитектура проекта уже подготовлена к дальнейшей модульной декомпозиции благодаря выделенным файлам:
    \begin{itemize}
        \item \texttt{selection.js};
        \item \texttt{resize-handles.js};
        \item \texttt{project-store.js};
        \item \texttt{html-exporter.js};
        \item \texttt{fld-serializer.js}.
    \end{itemize}
\end{itemize}

\subsection{Недостатки}

\begin{itemize}
    \item Основная логика приложения пока сосредоточена в файле \texttt{renderer.js}, из-за чего renderer-слой остаётся частично монолитным.

    \item Отсутствует полноценная адаптивная верстка экспортируемых страниц для мобильных устройств.

    \item Не реализована поддержка JavaScript-интерактивности экспортируемых компонентов.

    \item Нет импорта проектов из сторонних редакторов интерфейсов (например, Figma или Photoshop).

    \item Отсутствует совместная работа нескольких пользователей над одним проектом.

    \item Не реализована система undo/redo для отката изменений.

    \item Часть подготовленных модульных файлов пока содержит только заготовки и требует дальнейшего переноса логики из \texttt{renderer.js}.
\end{itemize}

\subsection{Выводы}

Применение языковой модели Codex в среде Visual Studio Code позволило существенно ускорить процесс разработки desktop-приложения на Electron.js. Модель использовалась как инструмент проектирования архитектуры, генерации исходного кода, исправления ошибок и подготовки вспомогательной документации.

В ходе работы Codex выполнял:
\begin{itemize}
    \item генерацию структуры проекта;
    \item создание HTML/CSS/JavaScript файлов;
    \item реализацию логики drag-and-drop;
    \item разработку системы хранения состояния;
    \item реализацию экспорта и сериализации проектов;
    \item исправление ошибок логических проблем.
\end{itemize}

Использование LLM позволило сократить время разработки и быстро получить рабочий прототип визуального конструктора интерфейсов. При этом основным ограничением остаётся необходимость проверки архитектурных решений и качества сгенерированного кода человеком, особенно при увеличении размера проекта и усложнении бизнес-логики.

\newpage
\section{Ссылка на исходный код}

Исходный код приложения и настоящий отчёт размещены в репозитории на GitHub:

\begin{center}
\url{https://github.com/sakuraruz/ui\_builder}
\end{center}

Структура репозитория:

\begin{itemize}
    \item \texttt{src/main/} --- главный процесс Electron-приложения; создание окна приложения, работа с IPC и экспортом проекта;

    \item \texttt{src/preload/} --- preload-слой, обеспечивающий безопасный мост между renderer-слоем и Node.js API;

    \item \texttt{src/renderer/} --- пользовательский интерфейс визуального конструктора: палитра элементов, холст, инспектор свойств, drag-and-drop и экспорт;

    \item \texttt{src/shared/} --- общие константы, структуры данных и вспомогательные типы;

    \item \texttt{docs/} --- документация по архитектуре проекта и внутреннему устройству приложения;

    \item \texttt{tests/} --- заготовки для автоматизированных тестов приложения;

    \item \texttt{src/renderer/index.html} --- основная HTML-разметка интерфейса редактора;

    \item \texttt{src/renderer/renderer.js} --- главный renderer-модуль приложения, содержащий основную бизнес-логику конструктора;

    \item \texttt{src/preload/index.js} --- preload API для взаимодействия renderer-слоя с Electron main-process;

    \item \texttt{src/main/index.js} --- точка входа Electron-приложения;

    \item \texttt{base.css} --- базовые стили приложения;

    \item \texttt{layout.css} --- стили layout-компонентов интерфейса;

    \item \texttt{canvas.css} --- стили рабочего холста и элементов редактора;

    \item \texttt{package.json} --- конфигурация Node.js/Electron проекта и зависимости;

    \item \texttt{README.md} --- техническое задание и описание проекта;

    \item \texttt{report.tex} --- исходный код настоящего отчёта в формате LaTeX.
\end{itemize}

\end{document}
